\honest{Infinite Certainty}{Eliezer Yudkowsky}{2008}

In ``Absolute Authority,'' yesterday, I argued that you do not need infinite certainty in order to choose.

A statement, such as the lightspeed limit, can be true always and exactly while our confidence in it falls short. If you flip a coin and do not look, it may be showing heads while you are completely unsure. ``A degree of uncertainty is not the same as a degree of truth or a frequency of occurrence.''

The same goes for arithmetic. I am not sure what ``true'' means for ``2 + 2 = 4,'' but I guess that in whatever sense it is true, it is exactly true, not true a trillion times minus one; that guess is more credible than any philosophical theory of what ``true'' means. Still, I am not absolutely confident of it. I could have hallucinated the evidence, or misremembered it. \nb{Standard probability theory gives probability 1 to every logical truth, and that SS0 + SS0 = SSSS0 follows from the Peano axioms is one. Some Bayesians, such as Hacking and Garber, relax this for real reasoners. The post takes their side without saying that the standard theory requires probability 1 here.}

What you should aim for is calibration: statements you hold at 99\% should come true 99 times in 100. That is harder than it sounds: ask a hundred people for ten statements each at 99\%, and do you think about ten of the thousand will be wrong? I will not describe the experiments on calibration, because people use them against opinions they dislike and forget them when they consider their own; they are in my book chapter. I only report that people who say they are 99\% sure are not right 99\% of the time. \nb{The chapter supports this. In one study it quotes, answers given odds of 100 to 1 were right 73\% of the time, and answers given a million to one or more were right 90\% of the time.}

To be 99.99\% sure that 2 + 2 = 4 is to claim that you could make 10,000 independent statements with equal confidence and be wrong about once. Maybe you could, for something so simple, both mathematical and empirical, and quietly taken for granted by everyone. I do not think you could for ``53 is a prime number'': to make 10,000 independent statements like it, each checked a new way, you would fail more than once. And a confidence of 99\% in 2 + 2 = 4 is a confidence that it is always and exactly true, not that it is true 99\% of the time. To claim 99.9999\% is to claim a million statements with about one error: a year of talking, at one statement every 20 seconds for 16 hours a day. \nb{``Absolute Authority'' had given the same illustration the day before.} A trillion statements would mean talking ``for a hundred human lifetimes'' without an error. \nb{At that pace a trillion statements take about 950,000 years, some twelve thousand lifetimes.} Claim 1 minus one over a googolplex, and your ego exceeds that of ``mental patients who think they're God.'' Even 1 minus one over 3 ↑↑↑ 3 is not much closer to 1 than 90\%. \nb{That depends on a scale the post does not give. It arrives the next day, in ``0 And 1 Are Not Probabilities.''}

If all else fails, the Dark Lords of the Matrix, tampering with your judgment of this very sentence, will stand in the way of infinite certainty. Am I absolutely sure of all this? ``Why, of course not.''

Last, Rafal Smigrodski: you can give less than certainty to the mathematics needed to derive Bayes's rule and still use it; and once you assign a probability of 1, ``I can never undo it,'' and must reject everything that disagrees. \nb{This is correct for updating by Bayes's rule. ``Absolute Authority'' had made the same point the day before.}
